\documentclass[pdflatex, sn-nature]{sn-jnl}
\usepackage{graphicx}
\usepackage{amsmath, amssymb}
\usepackage{booktabs}
\usepackage{longtable}
\newcommand{\widefigure}[1]{\makebox[\linewidth][c]{\includegraphics[width=160mm]{#1}}}
\raggedbottom
\begin{document}
\title[Affinity scores for target selection]{Local labels and measurement resolution shape the value of affinity scores}
\author{\fnm{Anonymous} \sur{Authors}}
\abstract{Computed interaction scores can rank proteins for a molecule, but benchmark rankings do not reveal what they add to locally available experimental information. We examine classical and learned outputs on a secondary-pharmacology panel and a kinase affinity panel, retaining quantitative measurements and assay bounds. Native binding-probability outputs exceed affinity outputs overall, yet their advantage contracts for pairs of quantitatively resolved interactions. On secondary pharmacology, Nesso augmentation raises first-test best-target recovery from 23.4\% to 35.6\% over chemistry at a quarter training budget, on 205 certifiable molecules; native binder scoring remains competitive without local fitting. Scaffold-held-out predictions for 47 kinase compounds attain 83.9\% target-order concordance using four representation components and 82.9\% using four scalar member outputs, compared with 77.4\% for native binder. Transfer to locally unlabelled proteins preserves useful native information, while augmentation over native scoring is smaller than improvement over a chemistry--sequence reference. Quantitative conclusions also depend on whether molecules or comparisons receive equal weight. Compact downstream models can therefore recover useful score information, but their practical value changes with the experimental decision and available local data.}
\maketitle

\section{Introduction}
A molecule rarely encounters only one protein. Its distribution of interactions helps determine which target explains an observed effect, which alternative target might support a new use, and which unintended interactions deserve experimental follow-up. Large biochemical profiling studies have made this problem measurable across kinase families and secondary-pharmacology panels~\cite{davis2011selectivity,sutherland2023spd}. Yet the computational question remains different from the familiar virtual-screening task. Choosing molecules for one protein requires comparisons along a column of an interaction matrix; choosing proteins for one molecule requires comparisons along a row. A score useful for one direction need not be equally useful for the other.

Docking and learned scoring functions provide rapid estimates of molecular compatibility, while recent foundation models combine biomolecular representations with structure and affinity prediction~\cite{trott2010vina,mcnutt2021gnina,passaro2025boltz2,shenoy2026nesso}. Their outputs carry different information. Boltz-2, for example, provides a binding-probability output intended for distinguishing binders from decoys and a continuous affinity output intended for comparing active compounds~\cite{boltzdocs}. Structural accuracy is another distinct property: FoldBench separates the ability to generate a good structure from the ability to select it~\cite{xu2025foldbench}. Progress in coordinate-robust affinity prediction and structure-aware generative models further broadens the available information~\cite{liu2026ecloudbind,cremer2026flowrroot}. A user deciding which target to test, however, needs to know which of these outputs improves that particular decision.

The experimental reference also shapes the question a benchmark answers. An assay may report a quantitative value for one interaction and only a bound for another. Correctly placing a measured interaction above a value beyond the assay's limit is useful, but does not establish accurate ordering among two measured interactions. Likewise, a molecule with many resolved targets supplies many more pairwise comparisons than a molecule with few. Pooling every comparison asks about a randomly selected comparison; weighting molecules equally asks about a randomly selected molecule. These are related but different experimental priorities.

Existing work has established that target priors, dataset composition and local adaptation can strongly affect predictive performance. TAPB studies target-prior bias in drug--target interaction prediction~\cite{lin2025tapb}. Concurrent work, MIRAGE, examines interpolation, redundancy and simple chemistry/protein reference models in affinity evaluation~\cite{yazdani2026mirage}. FLOWR.ROOT and recent Boltz-2 adaptation experiments show that local experimental labels can improve predictions in particular discovery settings~\cite{cremer2026flowrroot,furui2026boltzadaptation}. Molecular representation studies similarly find that additional modalities have task-dependent value~\cite{martinezgalindo2026multimodality,cui2026aca}. These findings motivate a more specific question: what information does a computed score add to within-molecule target selection beyond information already available from chemistry, protein sequence and local experiments?

Here we address that question with source-derived measurements, frozen native predictions and matched local-data comparisons. We distinguish native scoring from fitted calibration and augmentation; preserve assay bounds without creating an activity threshold; and vary local experimental information separately for new molecules and new proteins. The resulting comparison is an information audit as well as a leaderboard. It identifies useful positive cases, explains why aggregate score rankings can misrepresent a quantitative decision, and evaluates a practical remedy: fitting a modest readout of an existing score or representation alongside an inexpensive molecular reference model.

\begin{figure}[t]
\centering\widefigure{figs/graphical_abstract.pdf}
\caption{\textbf{Evaluate affinity scores for the molecular decision and available local data.} \textbf{a}, Target selection compares proteins for one molecule; the experimental NR3C1--dexamethasone structure 1M2Z illustrates a molecular interaction and is not a model prediction. \textbf{b}, Native Boltz binder-minus-affinity differences contract for two-exact comparisons. Whiskers are conditional 95\% paired molecule-bootstrap intervals. \textbf{c}, Chemistry and chemistry plus a fixed Nesso affinity output on SPD, across nested training-scaffold budgets. Lines show one fixed out-of-fold partition; its uncertainty is detailed in Fig.~\ref{fig:labels}.}
\label{fig:overview}
\end{figure}

\section{Results}
\subsection{Two measured panels expose different target-selection decisions}
The secondary-pharmacology dataset (SPD) supplies a chemically diverse drug panel and ten human target-associated binding-assay groups. The scored subset contains 930 molecules. An exact-identity join to the source records recovers 7,463 assay observations across 7,142 molecule--target cells; 1,308 records report quantitative IC$_{50}$ values and 6,155 report values above an assay limit. We analyze source-designated representative observations and a separate envelope across available protocols. These policies preserve the source's selection and measurement information rather than replacing bounds with exact values.

The ten groups represent KCNH2, ADRB1, SLC6A3, CHRM1, CHRM3, DRD2, SLC6A2, SLC6A4, NR3C1 and HRH1. They span 14 actual assay protocols because some transporter groups combine protocols. The endpoint is therefore a profile of protocol-specific binding-assay potency. For example, the KCNH2 assay measures ligand displacement, not ion-channel current. Its biological relevance motivates target prioritization, while the measured endpoint defines what our ranking analysis can establish.

DAVIS provides a complementary kinase setting: 72 compounds across 442 assay constructs, corresponding to 434 distinct recorded protein sequences and 31,824 measured cells~\cite{davis2011selectivity,wu2025daviscomplete}. Of these, 22,400 reach the source's 10\,$\mu$M K$_d$ reporting ceiling. Native analyses retain all 72 molecules and separately report 47 compounds with full standardized-InChIKey agreement between measurements and model inputs. All fitted comparisons use this 47-compound identity subset. Constructs with the same sequence string remain separate measurements and share target-held-out folds.

The SPD comparison includes 12 outputs from classical docking, learned rescoring, Boltz-2, supplied-pose Boltz affinity inference and Nesso. These are output channels and protocol variants, not 12 independent algorithms. DAVIS provides both Boltz outputs, their two ensemble members and 256 cached interaction-representation features. Native prediction acquisition is held fixed throughout the local-label experiments.

\subsection{Native output advantages contract when both interactions are quantitatively resolved}
On the representative SPD panel, 373 molecules supply 6,496 determinate comparisons; 261 molecules supply 2,235 two-exact comparisons. Native Boltz affinity achieves 56.2\% molecule-macro target-order concordance and its binder output achieves 65.9\%. The paired difference is 9.7 percentage points (conditional 95\% bootstrap interval 6.9 to 12.6). Among pairs with two exact measurements, concordance is 58.3\% and 61.9\%, respectively: a 3.7-point difference (interval $-0.6$ to 7.9). The corresponding difference for pairs containing one bounded observation is 11.2 points (Fig.~\ref{fig:heads}). Nesso also benefits from its binder output overall, but the difference is smaller among two-exact pairs.

DAVIS shows a larger contraction. Its native affinity and binder outputs achieve 69.3\% and 75.9\% overall, a 6.6-point advantage for the binder output (interval 4.9 to 8.4). With two exact measurements, concordance is 63.6\% and 63.7\%, and the paired difference is 0.1 points (interval $-1.8$ to 1.6). In the 47-compound identity subset, the two-exact estimates are 64.3\% and 63.1\%. Thus the aggregate binder advantage does not carry over unchanged to the quantitative comparison.

The ranking direction also matters. On SPD, native Boltz affinity attains 80.5\% ligand-order concordance but 56.2\% target-order concordance. Among two-exact pairs, these values are 63.4\% and 58.3\%. The large aggregate difference between ranking directions therefore combines a change in the decision with a change in its observable comparison mix. Table~\ref{tab:native} presents the complete SPD leaderboard. The strata describe naturally different observations; this is a decomposition of observed performance, not an intervention on experimental resolution.

\begin{figure}[t]
\centering\widefigure{figs/native_head_comparison.pdf}
\caption{\textbf{The native binder advantage contracts among quantitatively resolved interactions.} Paired differences in molecule-macro target-order concordance between Boltz-2 binder and affinity outputs on SPD (representative policy) and DAVIS (all 72 compounds). Points are observed differences; whiskers are 2.5 th--97.5 th percentiles from 5,000 paired molecule bootstrap samples. Each stratum retains the pairs whose order is established by the source measurements. Intervals condition on fixed predictions and target panels.}
\label{fig:heads}
\end{figure}
\begin{table}[t]
\caption{\textbf{Native SPD output performance depends on the comparison.} Percent concordance under the source-representative policy. Ligand order averages over targets; target order averages over molecules. All methods use common finite-score support. ``Exact'' requires two quantitative measurements. No local labels are fitted. Boltzina denotes supplied-pose affinity inference; the two missing-atom treatments are reported separately.}
\label{tab:native}
\centering\small\begin{tabular}{lrrr}
\toprule
Output & Ligand order & Target order & Target order, exact \\
\midrule
Boltz-2 affinity & 80.5 & 56.2 & 58.3 \\
Boltz-2 binder & 79.0 & 65.9 & 61.9 \\
Boltzina far affinity & 80.6 & 56.3 & 57.5 \\
Boltzina far binder & 79.6 & 67.0 & 63.7 \\
Boltzina zero affinity & 80.7 & 56.3 & 57.5 \\
Boltzina zero binder & 79.6 & 67.0 & 63.7 \\
GNINA on Vina pose & 69.2 & 61.3 & 58.8 \\
Nesso affinity & 75.5 & 55.5 & 54.2 \\
Nesso binder & 77.8 & 61.9 & 58.8 \\
AutoDock4 & 67.1 & 54.3 & 50.7 \\
GNINA own search & 70.3 & 62.8 & 60.5 \\
Vina & 64.7 & 58.0 & 49.5 \\
\bottomrule
\end{tabular}

\end{table}

\subsection{The added value of a score changes as local labels accumulate}
We next compare each score against a chemistry model trained on exactly the same local measurements. Molecules sharing a stereochemistry-preserving scaffold belong to the same outer and inner folds. Within each outer training partition we use nested fractions of one eighth, one quarter, one half and all eligible scaffold groups. Three inner folds select regularization without accessing the outer test measurements. The quarter budget corresponds to 138 training scaffold groups, 165--195 molecules and 1,230--1,509 measured cells per outer fit; the full budget contains 552 groups, 744 molecules and 5,505--5,618 measured cells. Thus these budgets describe existing local profiling information, not a few-shot assay claim. The chemistry reference combines a chiral Morgan-fingerprint kernel with six molecular descriptors; known targets have independent fitted coefficients. Both the reference and score-augmented models minimize the same loss to experimental intervals.

The chemistry model's representative-policy concordance rises from 63.5\% at the smallest budget to 71.1\% at the full budget. Adding Nesso affinity gives 66.6\%, 69.7\%, 71.6\% and 72.2\% across the four budgets. Its gain over chemistry is 3.9 percentage points at the quarter budget and 1.1 points at the full budget (Fig.~\ref{fig:labels}). The assay-envelope policy gives closely similar overall gains of 4.2 and 1.0 points. Matched availability-only controls do not explain this effect: the corresponding representative-policy increments over those controls are 4.0 and 1.3 points.

The same score has a different budget dependence for two-exact comparisons. Nesso affinity adds 0.5 points at the quarter budget and 2.0 points at the full budget under the representative policy; the envelope estimates are $-0.7$ and 2.0 points. Boltz binder augmentation helps the overall task but gives lower two-exact point estimates than chemistry across the displayed budgets. The full set of output-specific differences and their uncertainty is shown rather than reducing these curves to a single winning method.

Calibration alone is insufficient to explain the strongest fitted results. At the full representative-policy budget, a Boltz affinity-only calibration achieves 65.8\% and a Boltz binder-only calibration 65.5\%, compared with 71.1\% for chemistry. Nesso affinity-only calibration reaches 67.0\%, also below the chemistry reference and the 72.2\% augmentation result. The useful remedy is therefore not simply rescaling a native output. It is learning whether that output contributes information beyond the molecular reference under the available experimental budget.

\begin{figure}[p]
\centering\widefigure{figs/label_value_source_representative.pdf}
\caption{\textbf{Local labels change the incremental value of native scores.} SPD representative-policy comparisons: 373 evaluable molecules and 6,496 pairs overall; 261 molecules and 2,235 two-exact pairs. \textbf{a, b}, Chemistry and chemistry augmented by a single output across nested scaffold-training fractions, evaluated on all determinable and two-exact target pairs. \textbf{c, d}, Every native output added separately to chemistry at quarter and full budgets. Blue denotes all determinable pairs; orange denotes two-exact pairs. Bands and whiskers are conditional 95\% percentile intervals from 2,000 paired scaffold-group bootstrap samples of fixed five-fold predictions. They describe query sampling, without refitting or inference-seed variation. The corresponding envelope-policy figure and availability controls are provided in the supplementary material and Source Data.}
\label{fig:labels}
\end{figure}

\subsection{Target-relative score variation contains information beyond a molecule-wide summary}
A score may help a fitted model because it characterizes a generally high-scoring molecule, because it distinguishes that molecule's targets, or both. We separate each score profile into its mean across available targets and deviations from that mean. Mean-only and mean-plus-relative readouts use the same chemistry reference, local labels and original missingness indicator. The mean has target-dependent fitted coefficients, so this is a predictive comparison rather than the trivial observation that a constant cannot rank targets.

At the quarter representative-policy budget, the Nesso-affinity mean-only arm reaches 63.6\% overall concordance, below chemistry's 65.8\%. Adding target-relative variation raises it to 68.8\%, a 5.2-point increase (conditional 95\% interval 2.5 to 7.9). The envelope-policy increment is 6.6 points (4.1 to 9.1). For Boltz binder the corresponding increments are 2.9 and 3.8 points. The full-budget increments are smaller and less precisely resolved (Fig.~\ref{fig:components}). A molecule-wide summary therefore does not reproduce the useful low-budget profile information in these comparisons.

The decomposition is informative without being a uniformly better remedy. Raw-score augmentation reaches 69.7\% for Nesso affinity at the quarter representative budget, above the 68.8\% mean-plus-relative readout. Separately scaled components alter regularization and add flexibility; their purpose here is to locate predictive information, not establish a superior architecture. Moreover, both components require the original profile of native scores, so the mean-only arm does not by itself save acquisition cost.

\begin{figure}[p]
\centering\widefigure{figs/score_component_relative.pdf}
\caption{\textbf{Retaining target-relative variation improves the smaller-budget readout.} Change in target-order concordance when adding the relative component of a score profile to its molecule-wide mean, conditional on the same chemistry reference and original missingness indicator. Columns show Nesso affinity and Boltz binder; rows show all determinable and two-exact comparisons. Both assay policies and all four nested training budgets are retained. Whiskers are conditional 95\% paired scaffold-bootstrap intervals. The companion mean-versus-availability comparison and every output-specific contrast are provided in the supplementary material and Source Data.}
\label{fig:components}
\end{figure}

\subsection{Compact readouts recover kinase target profiles for new molecules}
The new-molecule DAVIS comparison gives a complementary positive result (Fig.~\ref{fig:readouts}). The smallest budget uses five locally characterized compounds and 2,210 measured cells per outer fold; the full budget uses 37--38 compounds and 16,354--16,796 measured cells. At the full local training budget, a readout of four principal components of the cached interaction representation achieves 83.9\% target-order concordance, compared with 77.4\% for the native Boltz binder output and 72.4\% for chemistry. The full 256-feature readout reaches 82.8\%, and the four scalar ensemble-member outputs reach 82.9\%. The four-component advantage over four scalar members is 1.0 percentage point (conditional 95\% interval 0.4 to 1.5). Thus useful target-profile information can be recovered with a compact fitted readout, while a comparison with chemistry alone would overstate the unique contribution of the internal representation.

This benefit also appears among two-exact comparisons. The four-component readout reaches 68.6\%, compared with 63.1\% for native binder and 58.9\% for chemistry; the four-scalar readout reaches 67.9\%. At the smallest local budget, four-component and four-scalar readouts both give 78.9\% overall concordance, whereas the full representation gives 72.5\%. These results favour controlling readout capacity when few locally characterized molecules are available. They also show that the shrinking SPD augmentation gain is not a universal label-budget law: the measured benefit depends on the panel and the local generalization problem.

\begin{figure}[p]
\centering\widefigure{figs/cached_readouts.pdf}
\caption{\textbf{Compact readouts recover kinase target profiles from cached predictions.} Five-fold scaffold-held-out DAVIS comparisons on 47 molecules with full identity agreement. \textbf{a, b}, Overall and two-exact concordance of chemistry, the native binder output and two matched four-dimensional readouts. The native curve is a fixed reference without local fitting. \textbf{c, d}, Paired difference between the four-component representation readout and the four scalar member outputs. Blue bands and whiskers are conditional 95\% intervals from 2,000 paired bootstrap samples; molecule and scaffold resampling coincide because each molecule has a distinct exact scaffold. Principal components and scaling are learned within each training partition.}
\label{fig:readouts}
\end{figure}

\subsection{Transfer to locally unlabelled targets retains useful information}
For DAVIS target transfer, we hold out protein-sequence groups while retaining measurements for the same 47 molecules on the remaining targets. A chemistry--protein reference combines molecular features with cached ESM-2 sequence embeddings. Predictions are compared only between targets held out together by the same fitted model. At the smallest training budget, which supplies 44 training sequence groups, adding the Boltz binder output improves molecule-macro concordance by 8.5 points (95\% interval 5.3 to 11.8). Reading the cached interaction representation gives a 9.0-point gain (6.3 to 11.8). With all eligible training groups, the chemistry--protein reference reaches 81.1\%; binder augmentation reaches 84.4\% and representation plus chemistry reaches 84.5\%.

The quantitative subset again gives a different answer (Fig.~\ref{fig:transfer}). At the full budget, the point gains in two-exact molecule-macro ordering are $-1.5$ points for either scalar output and $-0.8$ points for representation plus chemistry. This does not imply that the same models universally worsen quantitative ranking. On the identical comparisons, pooling pairs gives positive increments of 1.6, 1.6 and 2.8 points, respectively. A molecule with only one exact comparison and a molecule with many exact comparisons have equal weight in the first statistic and very different weights in the second. Restricting to a common query set changes the overall-versus-exact gain contrast by only 0.1--0.5 points; the weighting question is consequential in its own right.

The exhaustive SPD comparison withholds each of 45 target pairs while fitting on the remaining eight targets. Across the fixed panel, the chemistry--sequence reference reaches 51.2\% representative-policy concordance, native Boltz binder 65.9\%, and binder augmentation 66.8\%. Augmentation therefore gains 15.6 points over the reference but only 0.9 over the native score (95\% interval $-0.8$ to 2.6); the envelope-policy native-score increment is 0.6 points ($-1.1$ to 2.4). This is substantial native predictive information, not a 16-point calibration benefit. A different positive case is Vina: augmentation improves its two-exact ordering over native Vina by 6.5 points (0.8 to 12.1) under the representative policy and 6.7 points (0.9 to 12.4) under the envelope policy. The complete output-specific leaderboard and both comparison anchors are in the supplementary material.

The representation result also separates information availability from the quality of a final readout. At low label budgets, useful ordering can be recovered from cached internal features without rerunning the structural model. Its added benefit over a scalar readout is not uniform over budgets or quantitative endpoints. Local target withholding establishes transfer relative to these fitted readouts, not absence of the molecules or proteins from pretraining.

\begin{figure}[p]
\centering\widefigure{figs/target_transfer.pdf}
\caption{\textbf{Useful target-transfer information need not improve every averaging objective.} DAVIS comparisons on 47 full-identity-matched molecules and 442 constructs grouped into 434 sequences. \textbf{a, b}, Five-fold target-transfer curves using all determinable or two-exact comparisons; targets are compared only when held out by the same model. \textbf{c}, Full-budget increments on the identical two-exact comparisons, with equal molecule or equal comparison weight. Bands and whiskers are conditional 95\% percentile intervals from 5,000 paired molecule bootstrap samples. The change of sign reflects different weighting of observed comparisons, without changing predictions or excluding poorly performing molecules.}
\label{fig:transfer}
\end{figure}

\subsection{Recovering the best measured target is a separate decision}
A pairwise ordering statistic need not track recovery of the best target at a limited testing budget. We therefore evaluate every budget from one to ten targets per molecule. A target is certified best only when its reported lower potency bound is at least the upper bound of every other target. This rule uses the supplied measurements; it adds no active/inactive threshold. Among 909 molecules with all native scores available, the representative policy certifies a best target for 205 and the envelope policy for 198.

At the quarter representative-policy training budget, adding Nesso affinity raises recovery on the first test from 23.4\% to 35.6\%, a 12.2-percentage-point increase (conditional 95\% interval 4.0 to 20.5). The mean rank of the first certified best target improves from 3.76 to 3.20, a reduction of 0.56 ranks (interval 0.17 to 0.95). At the full training budget, the corresponding rank difference is only 0.08 in favour of augmentation (interval ranging from 0.25 worse to 0.43 better). Thus the positive effect translates to a concrete target-testing decision at the smaller budget, while the full curves retain performance across every testing budget and both measurement policies (Fig.~\ref{fig:decisions}).

Native binder scoring also remains a useful positive case: without fitting local labels, Boltz-2 recovers a certified best target on the first test for 39.5\% of these molecules and places it at mean rank 3.01. It is competitive with fitted alternatives for this short-list decision, despite lower overall pair-order concordance than some augmented models.

These ranks summarize a retrospectively identifiable subset. Observing one assay result in a new campaign would not itself reveal that the panel's best target had been found. The unresolved profiles remain in the release, with bounds on recovery rather than an invented best-target label.

\begin{figure}[p]
\centering\widefigure{figs/decision_curves.pdf}
\caption{\textbf{A useful pair-order score need not give the best short list.} Fraction of molecules whose first $k$ predicted targets include a certified best target, across every $k=1,\ldots, 10$. Exact prediction ties are averaged uniformly. Rows retain the representative and assay-envelope policies separately; columns compare quarter and full training budgets. Solid curves are fitted out-of-fold predictions; dashed curves are static native outputs. The grey band is a conditional 95\% scaffold-bootstrap interval for chemistry. All method-specific intervals and paired differences are supplied in Source Data. The evaluated subset contains 205 or 198 of 909 score-complete molecules; unresolved profiles are not labelled as failures or successes.}
\label{fig:decisions}
\end{figure}

\clearpage
\section{Discussion}
The practical value of an affinity score is conditional on the comparison a scientist needs and the information already available. Across these panels, a strong native binder output can be useful for broad target prioritization without having the same advantage for ordering quantitative interactions. Local labels recover additional useful information. In SPD and DAVIS target transfer, score increments are larger at smaller tested budgets. For new molecules on the fixed kinase panel, however, compact readouts retain substantial gains at the full budget. The difference concerns both the measured panel and the generalization problem. These findings preserve positive uses of modern binding models while giving a more precise interpretation to their aggregate rankings.

The distinction between native and fitted performance is central. A score-only calibration asks whether a local mapping can repair its scale; augmentation asks whether it supplies information beyond chemistry; a representation readout asks what can be recovered before the final scalar projection. They have different training and acquisition costs. In the completed DAVIS target-transfer block, the 200 fits for binder augmentation required 10.3 seconds of summed fitting wall time, compared with 784.3 seconds for representation plus chemistry. These are local post-processing costs with native predictions already cached, not the cost of obtaining those predictions. The current acquisition records span different devices, molecular subsets and preparation protocols, so they do not support a single comparable dollar-per-prediction leaderboard.

Training objectives provide a plausible explanation to test. Both Boltz-2 and Nesso combine absolute affinity supervision with a more heavily weighted term for relative differences within an assay~\cite{passaro2025boltz2,shenoy2026nesso}. That relative term directly supports comparing molecules under a common protocol, whereas target selection requires comparisons across proteins and often across protocols. The absolute term still constrains cross-assay predictions. Our readout experiments identify recoverable information; attributing the observed differences specifically to training losses would require an intervention on those losses.

Measurement resolution and averaging are equally practical choices. If the aim is to improve target choice for each incoming molecule, molecule-macro performance is directly relevant. If the aim is to maximize correctly ordered comparisons over a large measured collection, pooled performance answers that objective. The DAVIS sign change shows why a reporting convention cannot substitute for stating the scientific decision. Nor should a model selected using all determinable pairs be assumed optimal for a two-exact task: the present comparison holds selection fixed to isolate a reproducible operating procedure, and task-specific tuning remains a separate intervention.

The added knowledge is a quantitative map of which decision-relevant information remains in a score after local data are used. In the smaller SPD training regime, Nesso's score supplies a measurable first-test recovery gain over chemistry; in the same experiment, native Boltz binder scoring remains competitive without local fitting. As labels accumulate, the difference in mean first-best rank between chemistry and Nesso augmentation contracts. Coupled with the measurement-resolution and weighting results, this explains how several apparently conflicting score rankings can each be correct for a different experimental objective. It complements studies of pooled affinity generalization and target-specific virtual-screening adaptation~\cite{yazdani2026mirage,furui2026boltzadaptation}.

The scope is defined by the measured panels and preserved inference caches. SPD potency values depend on assay protocols, and its scored subset is smaller than the original resource. DAVIS contains few compounds, modified assay constructs and substantial ceiling censoring. Neither local holdout regime proves pretrained novelty. Original native inference is not equally replayable for every producer, although the released arrays, explicit axes and source hashes support complete reproduction of the downstream analyses. These limits motivate reporting the evaluated domain clearly, rather than converting a panel comparison into a universal ranking of algorithms.

For practice, the results suggest a concrete sequence. First specify whether the next experiment seeks a broad interaction, a quantitative distinction or the best target in a panel. Then compare each candidate output with a local chemistry/protein reference on the same labels, missingness and held-out units. A small readout of cached scores is a useful starting remedy when the observed gain justifies acquisition. Internal representations become worthwhile when they add value beyond those scalars for the chosen decision. The resulting choice can differ across projects without contradicting the underlying model's capabilities.

\section{Methods}
\subsection{Source measurements and identity joins}
SPD measurements, assay metadata and molecular structures derive from the published secondary-pharmacology resource~\cite{sutherland2023spd}. We join cached predictions by full InChIKey and retain source rows marked as exact structural matches. Only human protein-binding assay groups corresponding to the ten scored genes enter this analysis. Reported concentrations in micromolar are transformed as $pIC_{50}=6-\log_{10}(IC_{50})$. Exact observations become singleton intervals. A reported $IC_{50}>c$ becomes an upper-open interval $(-\infty, 6-\log_{10}c)$. Missing measurements remain missing.

The representative policy retains the source-designated record for each drug--assay-group pair. The envelope policy takes the smallest lower endpoint and largest upper endpoint across the matched protocols. Its upper boundary is open only if all records attaining that boundary are open. This envelope describes sensitivity to available assay protocols, not uncertainty from repeated raw measurements. The source has already summarized repeated observations, including its own handling of exact and censored values. No new activity or selectivity cutoff is introduced.

DAVIS records use the complete 442-construct table with molecular and protein identifiers. We align predictions through compound identifiers and construct names, rejecting duplicate keys. Reported nanomolar K$_d$ values become $pK_d=9-\log_{10}(K_d)$. Source values at 10,000 nM represent the reporting ceiling and become $(-\infty, 5]$. Full standardized-InChIKey matches define the 47-molecule fitting set independently of performance. The 72-molecule native comparison and the other identity categories remain available as sensitivity evidence. Identical recorded sequences share folds, including distinct assay constructs whose measurements differ.

\subsection{Native predictions and protein representations}
All native values are oriented so larger scores rank stronger predicted interactions. Docking energies and the relevant continuous affinity fields are negated; probabilities and CNNaffinity values retain their orientation. Their numerical values are not universally converted into experimental pK units. All native algorithms use their archived preparations and output-selection rules. The public acquisition record specifies verified versions, model revisions, sampling settings, masks and available source hashes.

Boltz-2 cofolding and supplied-pose Boltz inference use different preparation pathways; their contrast therefore does not isolate pose generation. GNINA on the first Vina pose is separate from GNINA's own search. Nesso affinity and binder outputs share a producer, as do the two Boltz heads. For DAVIS, the two 128-dimensional cached post-aggregation representations are concatenated. Four scalar member outputs provide a dimensionality control; a four-component representation readout provides a second control for new-molecule prediction, with principal components fitted only on the training partition.

Protein features for local target transfer are cached ESM-2 650M embeddings. DAVIS sequence hashes align 442 constructs to 434 unique source sequences. SPD uses full canonical sequences for eight proteins and the recorded KCNH2 and NR3C1 domains for the other two. Each fitting partition estimates its own feature centering and scale. Input sequence hashes and recorded model revisions are released; local withholding and pretrained exposure are separate properties.

\subsection{Interval-defined ranking and error}
Let $[\ell_{it}, u_{it}]$ denote the observed interval for molecule $i$ and target $t$, with open boundaries recorded separately. The order $t\succ s$ is determined when every value permitted for $t$ is larger than every value permitted for $s$. Equal exact values and overlapping intervals do not define a strict order. We score only such determined comparisons and give a prediction tie half credit. For the set $P_i$ of comparable target pairs,
\begin{equation}
 C_i=\frac{1}{|P_i|}\sum_{(t, s)\in P_i}\left[\mathbf{1}\{\widehat y_{it}>\widehat y_{is}\}+\frac{1}{2}\mathbf{1}\{\widehat y_{it}=\widehat y_{is}\}\right],
\end{equation}
where each pair is oriented by the observed order. Molecule-macro concordance averages $C_i$ over evaluable molecules. Pooled concordance sums credits and pair counts before division. Ligand-order analyses apply the corresponding calculation within each target. We report all determined comparisons, two exact observations, one bounded observation and two bounded observations separately. Comparisons use common finite-native support across the outputs in a dataset, including all fitted arms and their controls.

For fitted predictions, interval error is the distance to the observed interval; exact-only error is reported separately. An interval loss of zero for a bounded observation means consistency with that bound, not a known prediction error of zero. No error in experimental affinity units is assigned to an uncalibrated native score. Models that cannot supply the complete required prediction support are marked unavailable rather than evaluated on an easier subset.

\subsection{Matched local models and nested validation}
Molecular features comprise a 2048-bit chiral Morgan fingerprint of radius 3 and molecular weight, calculated logP, topological polar surface area, hydrogen-bond donor and acceptor counts, and rotatable-bond count. The molecular kernel is the equal-weight sum of fingerprint Tanimoto similarity and a radial-basis kernel on training-standardized descriptors, with fixed inverse dimension bandwidth. An exact numerical kernel basis and constant component are used. Known targets have an identity basis. In target transfer, the protein kernel is a normalized linear kernel on training-standardized ESM embeddings plus a constant component.

The fitted predictor combines molecular and target features with optional score or interaction-representation features. All arms minimize
\begin{equation}
 \frac{1}{|\Omega|}\sum_{(i, t)\in\Omega} d\!\left(\widehat y_{it},[\ell_{it}, u_{it}]\right)^2+\lambda\|\theta\|_2^2,
 \label{eq:objective}
\end{equation}
where $\Omega$ contains observed training cells, $d$ is distance to the interval, and the global intercept is unpenalized. Technical numerical conditioning rescales optimization coordinates without changing this objective. Fits must satisfy solver success checks; difficult fits use an exact active-set Newton fallback or continued optimization of the same objective. No scientific performance criterion determines numerical completion.

Five scaffold-group outer folds evaluate new molecules on known targets. Exact Murcko scaffold strings preserve stereochemistry; acyclic compounds are grouped by full identity. Within each outer fold, nested fractions $1/8,1/4,1/2,1$ of training groups determine the available labels. Three grouped inner folds select among ridge candidates $0.001,0.03,1$, divided by the number of training targets. Selection maximizes inner molecule-macro concordance on all determined pairs; ties select the larger ridge. A lack of evaluable inner comparisons is recorded as unavailable. Feature scaling, kernel bases and principal components are fitted within the current training partition.

For new proteins with known molecules, five outer folds group identical sequence hashes. Training budgets select nested fractions of eligible protein groups, and inner folds also withhold protein groups. Pairwise evaluation compares only targets predicted in the same held-out context. It never ranks scores stitched together from separately fitted target-fold models. SPD additionally permits every unordered pair of its ten targets to be withheld together while fitting on the remaining eight; this exhaustive comparison is recorded separately from the five-fold partition.

A score-only calibration uses the score and its missingness indicator with target-dependent coefficients. Augmentation adds the same pair features to the chemistry/protein reference. An availability-only control supplies the same mask without the score. Identical masks share a control fit. Representations are training-standardized, with missing indicators retained where applicable. Individual outputs are added separately; no ensemble of all scoring methods is used to define the primary comparison.

\subsection{Information decomposition and decision curves}
For a fixed SPD score matrix $S$, the molecule-wide component is $M_i=\operatorname{mean}_{t: S_{it}\text{ observed}}S_{it}$ and the target-relative component is $R_{it}=S_{it}-M_i$. Both retain the original availability mask. Mean-only and mean-plus-relative arms are added to the same chemistry reference, with one original missingness indicator and otherwise unchanged folds, training budgets, candidate regularization and interval objective. Their comparison addresses predictive information in a fixed panel. Obtaining $M_i$ still requires the available target-profile scores; this decomposition alone implies no inference saving.

For best-target recovery, a certified set $G_i$ contains targets whose finite lower bound is at least every other upper bound. Missing observations admit the full interval $(-\infty,+\infty)$. Prediction ties are averaged over uniform permutations within each tie block. The recovery curve reports the probability that the top $k$ predictions intersect $G_i$ for every $k=1,\ldots, 10$. Conditional curves use molecules with a nonempty certified set and scores for all targets and methods. For the remaining score-complete population, certified and undominated target sets provide lower and upper recovery bounds. These are measurement-identification bounds, not sampling confidence intervals.

\subsection{Uncertainty, computational cost and reproducibility}
Paired bootstrap samples reuse the same sampled units across methods and training budgets. Native and DAVIS target-transfer summaries resample molecules; SPD fitted comparisons and decision curves resample molecular scaffold groups. Reported 95\% intervals are percentile intervals conditional on the saved predictions, labels, target panel and partitions. They do not represent retraining, new target-panel or model-inference variation, and do not define a threshold of scientific importance. Every query contribution, denominator and unavailable case is retained.

Fitting times are measured wall times for the recorded local jobs and are separated from original native acquisition. Summing fit times describes the logged fitting workload, not CPU utilization or a prospective deployment price. Archived native acquisition records are attributed only to their verified producer and subset. We do not extrapolate an unmatched device/subset measurement to the complete scored panel.

The analysis checks input bytes against SHA-256 manifests, reconstructs source intervals, and saves fitted predictions, partitions, hyperparameter selections and numerical diagnostics. Independent numerical replay verifies pair credits and aggregate results from the saved arrays. No new paid GPU computation was required for the reported downstream analyses.

\section*{Data availability}
Source datasets are attributed to the SPD resource~\cite{sutherland2023spd} and DAVIS~\cite{davis2011selectivity,wu2025daviscomplete}. The accompanying repository contains the analysis inputs, explicit identifier axes, preserved prediction caches, source checksums and machine-readable figure data. Experimental structure 1M2Z is available from the Protein Data Bank~\cite{pdb1m2z}. The final repository address is supplied in the release version of this manuscript.

\section*{Code availability}
The repository provides scripts to rebuild measurements, rerun the specified local fits, reconstruct every metric and regenerate the figures and PDF. Reproduction from the released native caches is distinguished from rerunning all original native acquisition workflows.

\begin{thebibliography}{99}

\bibitem{davis2011selectivity} Davis, M. I. et al. Comprehensive analysis of kinase inhibitor selectivity. \textit{Nature Biotechnology} \textbf{29}, 1046-1051 (2011). \url{https://doi.org/10.1038/nbt.1990}.

\bibitem{sutherland2023spd} Sutherland, J. J. et al. A preclinical secondary pharmacology resource illuminates target-adverse drug reaction associations of marketed drugs. \textit{Nature Communications} \textbf{14}, 4323 (2023). \url{https://doi.org/10.1038/s41467-023-40064-9}.

\bibitem{trott2010vina} Trott, O. \& Olson, A. J. AutoDock Vina: Improving the speed and accuracy of docking with a new scoring function, efficient optimization, and multithreading. \textit{Journal of Computational Chemistry} \textbf{31}, 455-461 (2010). \url{https://doi.org/10.1002/jcc.21334}.

\bibitem{mcnutt2021gnina} McNutt, A. T. et al. GNINA 1.0: molecular docking with deep learning. \textit{Journal of Cheminformatics} \textbf{13}, 43 (2021). \url{https://doi.org/10.1186/s13321-021-00522-2}.

\bibitem{passaro2025boltz2} Passaro, S. et al. Boltz-2: Towards accurate and efficient binding affinity prediction. \textit{bioRxiv} (2025). \url{https://doi.org/10.1101/2025.06.14.659707}.

\bibitem{shenoy2026nesso} Shenoy, N. et al. Nesso-1: Accelerating open-source binding affinity predictions. (2026). \url{https://www.valencelabs.com/wp-content/uploads/2026/07/nesso1.pdf}.

\bibitem{boltzdocs} Boltz contributors. Boltz prediction documentation: affinity outputs. \url{https://github.com/jwohlwend/boltz/blob/main/docs/prediction.md} (accessed 4 October 2026).

\bibitem{xu2025foldbench} Xu, S. et al. Benchmarking all-atom biomolecular structure prediction with FoldBench. \textit{Nature Communications} \textbf{17}, 442 (2025). \url{https://doi.org/10.1038/s41467-025-67127-3}.

\bibitem{liu2026ecloudbind} Liu, Y. et al. An electron-density point-cloud framework for robust protein-ligand interaction prediction. \textit{Nature Communications} \textbf{17}, 7424 (2026). \url{https://doi.org/10.1038/s41467-026-74196-5}.

\bibitem{cremer2026flowrroot} Cremer, J. et al. FLOWR.ROOT -- A flow matching-based foundation model for joint multi-purpose structure-aware 3D ligand generation and affinity prediction. \textit{Nature Communications} \textbf{17}, 5883 (2026). \url{https://doi.org/10.1038/s41467-026-74130-9}.

\bibitem{lin2025tapb} Lin, G. et al. TAPB: an interventional debiasing framework for alleviating target prior bias in drug-target interaction prediction. \textit{Nature Communications} \textbf{16}, 10867 (2025). \url{https://doi.org/10.1038/s41467-025-66915-1}.

\bibitem{yazdani2026mirage} Yazdani-Jahromi, M. et al. MIRAGE: Measuring Interpolation and Redundancy in Affinity GEneralization. (2026). \url{https://www.alphaxiv.org/pdf/2609.14491}.

\bibitem{furui2026boltzadaptation} Furui, K. \& Ohue, M. Adapting Boltz-2 with limited experimental activity data improves early enrichment in virtual screening. (2026). \url{https://www.alphaxiv.org/pdf/2609.24302}.

\bibitem{martinezgalindo2026multimodality} Martinez Galindo, M. et al. Impact of molecular multimodality on neural network models for prediction tasks related to drug discovery. \textit{Nature Communications} \textbf{17}, 8081 (2026). \url{https://doi.org/10.1038/s41467-026-74487-x}.

\bibitem{cui2026aca} Cui, C. et al. Activity-cliff awareness enables robust graph learning for molecular property prediction. \textit{Nature Communications} \textbf{17}, 9376 (2026). \url{https://doi.org/10.1038/s41467-026-75713-2}.

\bibitem{wu2025daviscomplete} Wu, M. H. et al. Towards Precision Protein-Ligand Affinity Prediction Benchmark: A Complete and Modification-Aware DAVIS Dataset. (2025). \url{https://www.alphaxiv.org/pdf/2512.00708}.

\bibitem{pdb1m2z} Bledsoe, R. K. et al. Crystal structure of the glucocorticoid receptor ligand binding domain reveals a novel mode of receptor dimerization and coactivator recognition. \textit{Cell} \textbf{110}, 93--105 (2002). Experimental coordinates: \url{https://doi.org/10.2210/pdb1M2Z/pdb}.

\end{thebibliography}

\end{document}
